\scanpage{62}\mspage{61}
\deleted{\unsure{Soit|soit}} Ainsi, les \unsure{systèmes|système} $(\mathcal F,\leq)$ satisfaisant $C1,C2,C3$ et \unsure{tels|tel} que toute $F\in\mathcal F$ soit ``finie'' (i.e. l'ens $\mathcal F_{\leq F}$ des sous-figures est fini, ou encore, l'ens $M_{\leq F}\inserted{=\unsure{\overline F|F}}$ \deleted{des} de \deleted{\unsure{sous}} ses strates est fini) \unsure{équivaut|équiv.} à :

$(M,\leq,R)$ \deleted{\unsure{str.}} où $\leq$ relation d'ordre, $R$ relation refl. et sym., satisfais. \unsure{à|a}
\[
\begin{aligned}
\text{a) }&\text{les }M_{\leq x}\text{ finis}\quad(x\in M)\\
\text{b) }&\text{Si }(x,y)\in R\text{ et }x'\leq x,y'\leq y,\ \text{\unsure{on a} }\deleted{\unsure{R}}(x',y')\in R\,;
\end{aligned}
\]
où les figures sont \unsure{interprétées|interprète|interprètent} comme les parties \uline{fermées}, \uline{finies} \inserted{$A$} de $M$, telles que $\forall x,y\in A$, on ait $(x,y)\in R$.

% slanted
\margin{\unsure{Ainsi} la notion \uline{\unsure{combinatoire|constructive|construction}} de \unsure{connexité|compacité}. Elle \illegible{1} \unsure{en évidence|coïncide} \deleted{\unsure{qui|que}} \unsure{avec|aussi} la notion \illegible{1} \unsure{qui la sous-entend}, \unsure{pour|que} les \unsure{str. top.|strates} sont 0-connexes + \ldots\ En fait, si \inserted{on \unsure{vient de|veut|voulait}} \deleted{``\unsure{exprimer|exiger|exigé}''} \unsure{$=$|--} que les \unsure{conditions} pr. $A$ \unsure{des|clos} \deleted{\unsure{et|de}} de $x$ \unsure{soient|soit} \unsure{co}}

\uline{Composantes connexes} Une figure est dit \uline{connexe} si \deleted{\unsure{on a}}
\[
F=F'\cup F''\quad
\inserted{\text{avec }F',F''\ \text{\uline{disjoints} i.e. }F'\cap F''=\phi}
\ \Longrightarrow\ F'=\deleted{\unsure{\phi}}F\ \text{ou }F''=\deleted{\unsure{\phi}}F.
\]
\margin{\uline{NB} $F$ connexe ss $\overline F$ l'est comme ens. \unsure{\deleted{filtré}|filtré} ordonné}

% slanted
\margin{\unsure{On écrit|Car ici} $F=F'\amalg F''$ (\unsure{signe de|Chacun de} \unsure{la somme|ces deux f.} \unsure{disjointe|disjoint}\unsure{.|)} (\uline{NB} Si \deleted{\unsure{deux|divise}} \unsure{parmi les|les} \deleted{\unsure{fig.}} $F',F''\unsure{,|\leq}F$ sont connexes, la \unsure{troisième} aussi.}

Une figure est dite \uline{0-connexe} si elle est connexe et $\ne\phi_{\mathcal F}$. Une sous-figure $F_0$ de $F$ est appelée \uline{composante connexe} de $F$ si a) elle est 0-connexe et

b) \deleted{$F_0$ est une \unsure{compos|comp.}} \inserted{$F=F_0\amalg \deleted{\unsure{F}}G$, $G$ comp.ble} \deleted{avec \unsure{$F$,|$F_1$} disjoint de $F_0$.}

\deleted{\unsure{Il est}} \deleted{\uline{\unsure{NB}}}\inserted{Lemme)} Soit $F=F'\amalg F''$, \deleted{\unsure{avec} $F'\cap F''=\phi$}, et soit $G$ une sous-figure connexe de $G$. Alors $G$ est contenue \unsure{ds|dans} $F'$ ou $F''$.
